\scanpage{140}
$\hookrightarrow$ extension sur $X$ invariante par translation
\[
0\mathbin{-}\unsure{\underline O_X|\underline t_X|\underline\omega_X}\mathbin{-}F\mathbin{-}\underline O_X\mathbin{-}0
\]
et un $G$-isom. $P\simeq\mathbb P(F)$ ($\simeq\overline{\mathbb V(\mathcal T)}$, où $\mathcal T$ est
le $G$-torseur sous $\underline O_X$, avec la $G$-structure
habituelle de $\underline O_X$, définie par l'extension \ldots)

[La \deleted{section}\inserted{horiz. $\sigma'$} de $P$ qui y correspond est
\unsure{partout|toute} disjointe de $\sigma$, donc correspond à
un splitting \inserted{(inv. par $H$)} $\mathscr J\simeq\underline\omega\oplus\underline O_X\oplus\underline\omega^{-1}$,
[qui d'ailleurs n'est autre que celui
déduit du splitting de Campbell-Hausdorff
sur l'alg. tronquée \ldots) avec $\underline O_X+\underline\omega^{-1}$ \inserted{et $\underline\omega^{-1}$} invar$^{\mathrm ts}$
pour la connexion ; et celle-ci induisant
la connexion \deleted{nulle} triviale sur le
quotient $\underline O_X$ \unsure{\ldots}] Pour les conn. affines
invariantes que l'on déduit par
la méthode de la ``conique'', il
y a $\hookrightarrow$ réduction \inserted{invariante} du groupe structural \deleted{au}
à un \inserted{un conjugué de,} \unsure{stabilis.} \unsure{sour|sous} de $G_0=G_t(\mathbb P(\underline\omega_0+\underline O_S))$.
Si on est sur un corps alg. clos (suffisant ici), \unsure{donc|alors}
la connexion proj. can. est caractérisée, parmi
celles qui sont invariantes, par les conditions
\deleted{de}\inserted{que} $P$ est \unsure{astreinte|admette} qu'une seule \deleted{\unsure{connexion|horizontale}}\inserted{section plate}
\deleted{\unsure{lin.}} invariante (\uline{NB} en car. $0$, \deleted{invariante}\inserted{plate}
? $\mid$ $\Rightarrow$ invariante), ttes les autres en
ont exactement deux \ldots.
% margin note written sideways
\margin{et si $T_0$ est le sous \unsure{tors.|tour} correspondant à $G_0$ (\unsure{uniqt déterminé|image déterminant}) on a $\unsure{\mathfrak P|\mathbb P}\simeq f^*(\overline{\mathcal T_0})=f^*(\overline{T_0})$ (\unsure{Algèbre projective|algébro-géomètre}), -- \unsure{conséquence} : c$^{\text{e}}$ section $1$, et la \deleted{\unsure{consid}}\unsure{restriction} \deleted{\unsure{elle}} à l'infini \deleted{\unsure{étant}} plate.}

\uline{NB} Si $H$ est propre sur $S$, fibres connexes (donc
$X/S$ un schéma en courbes elliptiques)
\uline{\unsure{tous|ttes}} les connexions projectives sur $X$
sont inv. par $H$. Ce qui caractérise donc
la conn. \inserted{proj.} canonique parmi \uline{t$^{\mathrm{tes}}$}
les connexions \ldots
% archive sheet number: 139
